% ch20_enriched_category_theory.tex — Volume II Chapter 8

\chapter{Enriched Category Theory}

\begin{overviewbox}
Chapter~12 introduced $\mathcal{V}$-enriched categories, where hom-sets are
replaced by objects of a monoidal category $\mathcal{V}$. This chapter develops
the deep theory: the enriched Yoneda lemma, enriched adjunctions, weighted
limits and colimits, base change, and self-enrichment of closed categories.

The enriched Yoneda lemma is the heart of the chapter. It states that the
``set of natural transformations from a representable $\mathcal{V}$-functor to
$F$'' is (up to the right notion of $\mathcal{V}$-natural transformation) just
$F(a)$ in $\mathcal{V}$ --- the same statement as ordinary Yoneda, but now
the conclusion lands in $\mathcal{V}$, not $\mathbf{Set}$.

The payoff: enriched limits, enriched adjunctions, and the general theory of
weighted limits that subsume ordinary limits as a special case. Weighted limits
are the correct notion of limit in the enriched world, and understanding them
is essential for working with chain complexes, spectra, and higher-categorical
structures.
\end{overviewbox}


% ═══════════════════════════════════════════════════════════════════════════
\begin{nameorigin}
\term{Enriched category theory} as a systematic discipline was
developed by \emph{Max Kelly} (1930--2007, Australian mathematician)
in his 1982 book \emph{Basic Concepts of Enriched Category Theory}.
The foundational papers by \emph{Eilenberg-Kelly} (1966) established
the theory. The slogan: \emph{ordinary category theory is the
$\mathbf{Set}$-enriched case of enriched category theory}, and many
constructions generalize uniformly across enrichments.

\term{Weighted limit} (sometimes \term{indexed limit}) was developed
to extend ordinary limits to the enriched setting, where the index
category is replaced by a $\mathcal{V}$-functor (the \emph{weight}).
The weighted limit $\{W, F\}$ generalizes ordinary limits ($W = $ const
$\mathbf{1}$) and many other constructions (cotensors, ends). Kelly's
book is the standard reference.
\end{nameorigin}

\section{The Enriched Yoneda Lemma}
\label{sec:enriched-yoneda}
% ═══════════════════════════════════════════════════════════════════════════

\subsection{Informally: Yoneda in $\mathcal{V}$}

The ordinary Yoneda lemma says: for a locally small category $\mathcal{C}$ and
a functor $F : \mathcal{C} \to \mathbf{Set}$, natural transformations from the
representable $\mathcal{C}(a, -)$ to $F$ are in bijection with elements of $F(a)$:
\begin{equation*}
  [\mathcal{C}, \mathbf{Set}](\mathcal{C}(a,-), F) \;\cong\; F(a).
\end{equation*}
When $\mathcal{C}$ is enriched over $\mathcal{V}$, the hom-objects $\mathcal{C}(a,b)$
are objects of $\mathcal{V}$, not sets. The Yoneda lemma must be rewritten to
have its conclusion in $\mathcal{V}$ rather than $\mathbf{Set}$, and
``natural transformations'' must be replaced by the \emph{$\mathcal{V}$-natural
transformations} that live in $\mathcal{V}$.

\begin{theorem}[$\mathcal{V}$-Yoneda Lemma]
\label{thm:enriched-yoneda}
Let $\mathcal{C}$ be a $\mathcal{V}$-enriched category and $F : \mathcal{C}
\to \mathcal{V}$ a $\mathcal{V}$-functor. For each object $a \in \mathcal{C}$,
there is a $\mathcal{V}$-natural isomorphism
\begin{equation*}
  [\mathcal{C}, \mathcal{V}]_\mathcal{V}(\mathcal{C}(a,-),\; F)
  \;\cong\;
  F(a),
\end{equation*}
where the left-hand side is the $\mathcal{V}$-object of $\mathcal{V}$-natural
transformations.
\end{theorem}

\subsection{Expanding: What the Statement Means}

The $\mathcal{V}$-object of $\mathcal{V}$-natural transformations between two
$\mathcal{V}$-functors $G, H : \mathcal{C} \to \mathcal{V}$ is defined as the
\emph{end}:
\begin{equation*}
  [\mathcal{C}, \mathcal{V}]_\mathcal{V}(G, H)
  \;=\;
  \int_{c \in \mathcal{C}} \mathcal{V}(G(c), H(c)).
\end{equation*}
This is an object of $\mathcal{V}$ built from all the individual hom-objects
$\mathcal{V}(G(c), H(c))$, glued together by the end construction (Chapter~26).

For the Yoneda case, $G = \mathcal{C}(a,-)$ and $H = F$:
\begin{equation*}
  [\mathcal{C}, \mathcal{V}]_\mathcal{V}(\mathcal{C}(a,-), F)
  \;=\;
  \int_{c} \mathcal{V}(\mathcal{C}(a,c), F(c)).
\end{equation*}
The Yoneda lemma says this end is isomorphic to $F(a)$ in $\mathcal{V}$.

\subsection{Formalizing: Full Proof via Prooflog}

\begin{prooflog}
\proofstep{We construct a morphism $\phi : \int_c \mathcal{V}(\mathcal{C}(a,c),
F(c)) \to F(a)$.}{We want the forward direction of the isomorphism.}
\proofstep{An element of the end $\int_c \mathcal{V}(\mathcal{C}(a,c), F(c))$
is a compatible family of morphisms $(\theta_c : \mathcal{C}(a,c) \to F(c))_c$
in $\mathcal{V}$.}{By the universal property of the end.}
\proofstep{Apply $\theta_a : \mathcal{C}(a,a) \to F(a)$ to the identity
$j_a : I \to \mathcal{C}(a,a)$ (the enriched identity map).}{The identity
morphism gives a specific element of $\mathcal{C}(a,a)$ over $I$.}
\proofstep{Define $\phi(\theta) = \theta_a \circ j_a : I \to F(a)$.}{This is
a morphism in $\mathcal{V}$.}
\proofstep{We construct the inverse $\psi : F(a) \to \int_c \mathcal{V}(\mathcal{C}(a,c), F(c))$.}{
For this we need to produce a compatible family $(\psi(x)_c)_c$ from $x : I \to F(a)$.}
\proofstep{At each $c$, define $\psi(x)_c$ as the composite $\mathcal{C}(a,c)
\xrightarrow{\cong} I \otimes \mathcal{C}(a,c) \xrightarrow{x \otimes \mathrm{id}}
F(a) \otimes \mathcal{C}(a,c) \xrightarrow{F_{a,c}} F(c)$.}{The action of $F$
on hom-objects is the $\mathcal{V}$-functor structure map $F_{a,c} : F(a) \otimes
\mathcal{C}(a,c) \to F(c)$.}
\proofstep{The family $(\psi(x)_c)_c$ is compatible (i.e., defines an element
of the end) because $F$ satisfies the $\mathcal{V}$-functoriality axioms.}{The
end condition is: for $f : c \to c'$ in $\mathcal{C}$, the square
$\mathcal{V}(\mathcal{C}(a,c'), F(c')) \to \mathcal{V}(\mathcal{C}(a,c),F(c))$
commutes. This follows from functoriality of $F$.}
\proofstep{$\phi \circ \psi = \mathrm{id}_{F(a)}$: compute $\phi(\psi(x)) =
\psi(x)_a \circ j_a$.}{Substituting the definitions.}
\proofstep{$\psi(x)_a \circ j_a = F_{a,a}(x \otimes j_a) \circ \lambda^{-1}_{\mathcal{C}(a,a)}
= F_{a,a}(x \otimes j_a) \circ \lambda^{-1} = x$.}{The $\mathcal{V}$-functor
axiom says $F_{a,a}(x \otimes j_a) = x$ when applied with the enriched unit
$j_a$ (unit axiom for $\mathcal{V}$-functors).}
\proofstep{$\psi \circ \phi = \mathrm{id}$: For a compatible family $(\theta_c)$,
we need $\psi(\phi(\theta))_c = \theta_c$ for all $c$.}{The converse direction.}
\proofstep{$\psi(\phi(\theta))_c = F_{a,c}((\theta_a \circ j_a) \otimes \mathrm{id}
_{\mathcal{C}(a,c)})$.}{Substituting $\phi(\theta) = \theta_a \circ j_a$.}
\proofstep{$= F_{a,c}(\theta_a \otimes \mathrm{id}) \circ (j_a \otimes \mathrm{id}) = \theta_c$}{
By the naturality (end condition) for $(\theta_c)$: the compatibility of $\theta$
with morphisms in $\mathcal{C}$ translates exactly to this equality.}
\proofstep{Both composites are identities; $\phi$ and $\psi$ are inverse
isomorphisms in $\mathcal{V}$.}{The enriched Yoneda isomorphism.}
\proofstep{$\mathcal{V}$-naturality in $a$: the isomorphism is natural in $a$
because all constructions (end, functor action) are natural.}{Standard argument:
the maps $\phi_a$ and $\psi_a$ assemble into $\mathcal{V}$-natural transformations.}
\end{prooflog}

\subsection{Reorganizing: Corollaries}

\begin{corollary}[Enriched Representability]
A $\mathcal{V}$-functor $F : \mathcal{C} \to \mathcal{V}$ is representable iff
there exists $a \in \mathcal{C}$ with $F \cong \mathcal{C}(a,-)$ as
$\mathcal{V}$-functors.
\end{corollary}

\begin{corollary}[Enriched Yoneda Embedding]
The enriched Yoneda embedding $Y : \mathcal{C} \to [\mathcal{C}^{\mathrm{op}},
\mathcal{V}]_\mathcal{V}$ is fully faithful as a $\mathcal{V}$-functor.
\end{corollary}

\begin{dialog}
Is the enriched Yoneda lemma really the same as the ordinary one?

Yes, in a precise sense. If $\mathcal{V} = \mathbf{Set}$ (with cartesian
product), then $\mathcal{V}$-enriched categories are ordinary locally small
categories, ends are ordinary natural transformation sets, and the enriched
Yoneda lemma reduces to the ordinary one. The enriched version is a strict
generalization.
\end{dialog}


% ═══════════════════════════════════════════════════════════════════════════
\section{$\mathcal{V}$-Adjunctions}
\label{sec:enriched-adjunctions}
% ═══════════════════════════════════════════════════════════════════════════

\subsection{Informally: Hom-Object Bijections}

In ordinary category theory, an adjunction is a bijection of hom-sets:
$\mathcal{D}(FA, B) \cong \mathcal{C}(A, GB)$. In the enriched world, hom-sets
are replaced by hom-objects in $\mathcal{V}$, and the bijection is replaced by
an isomorphism in $\mathcal{V}$:
\begin{equation*}
  \mathcal{D}(FA, B) \;\cong\; \mathcal{C}(A, GB)
  \quad \text{(isomorphism in } \mathcal{V}\text{)}
\end{equation*}
$\mathcal{V}$-natural in $A$ and $B$. This is a \term{$\mathcal{V}$-adjunction}
(or \term{enriched adjunction}).

\subsection{Expanding: Unit and Counit in the Enriched Case}

The unit of a $\mathcal{V}$-adjunction is a $\mathcal{V}$-natural transformation
$\eta : \mathrm{Id}_\mathcal{C} \Rightarrow GF$. But the naturality of $\eta$ is
now $\mathcal{V}$-naturality: the naturality square is a commutative diagram in
$\mathcal{V}$, not merely in $\mathbf{Set}$.

Explicitly, $\mathcal{V}$-naturality of $\eta$ at $A, A'$ is the commutativity
of the diagram in $\mathcal{V}$:
\begin{equation*}
\begin{tikzcd}[column sep=large]
  \mathcal{C}(A, A') \arrow[r, "G_{FA,FA'}F_{A,A'}"] \arrow[d, "\eta_{A'*}" '] &
  \mathcal{C}(GFA, GFA') \arrow[d, "\eta_A^*"] \\
  \mathcal{C}(A, GFA') \arrow[r, equal] & \mathcal{C}(A, GFA')
\end{tikzcd}
\end{equation*}
where $\eta_{A'*}$ and $\eta_A^*$ denote post- and pre-composition with $\eta$.

\subsection{Formalizing: The Definition}

\begin{definition}
Let $\mathcal{C}$ and $\mathcal{D}$ be $\mathcal{V}$-categories. A
\term{$\mathcal{V}$-adjunction} $F \dashv_\mathcal{V} G$ (with
$F : \mathcal{C} \to \mathcal{D}$ and $G : \mathcal{D} \to \mathcal{C}$
$\mathcal{V}$-functors) consists of a $\mathcal{V}$-natural isomorphism
\begin{equation*}
  \varphi_{A,B} : \mathcal{D}(FA, B) \;\xrightarrow{\;\cong\;}\; \mathcal{C}(A, GB)
\end{equation*}
for all $A \in \mathcal{C}$, $B \in \mathcal{D}$, natural in both $A$ and $B$
in the $\mathcal{V}$-sense.
\end{definition}

\begin{remark}
Every $\mathcal{V}$-adjunction has an underlying ordinary adjunction: applying
$\mathcal{V}(I, -)$ to the hom-object isomorphism gives an ordinary bijection
of hom-sets (since $\mathcal{V}(I, \mathcal{C}(A,B))$ is the set of morphisms
$A \to B$ in the underlying ordinary category of $\mathcal{C}$).
\end{remark}

\subsection{Reorganizing: Enriched vs.\ Ordinary Adjunctions}

The gap between ordinary and enriched adjunctions is subtle: an ordinary
adjunction between $\mathcal{V}$-categories (functors that happen to be
$\mathcal{V}$-functors, but with only a set bijection on hom-sets) need not
be a $\mathcal{V}$-adjunction. The stronger requirement of isomorphism in
$\mathcal{V}$ is genuinely richer.

However, in many practical cases (when $\mathcal{V}$ is the base of an enrichment
and the relevant categories are ``nice''), ordinary adjunctions between
$\mathcal{V}$-functors automatically lift to $\mathcal{V}$-adjunctions. This
is guaranteed, for example, when $\mathcal{C}$ and $\mathcal{D}$ are
$\mathcal{V}$-tensored and cotensored.


% ═══════════════════════════════════════════════════════════════════════════
\section{Weighted Limits and Colimits}
\label{sec:weighted-limits}
% ═══════════════════════════════════════════════════════════════════════════

\subsection{Informally: Why Ordinary Limits Are Not Enough}

In ordinary category theory, a limit of $F : \mathcal{J} \to \mathcal{C}$ is
characterized by $\mathcal{C}(c, \lim F) \cong [\mathcal{J}, \mathbf{Set}](\Delta 1,
\mathcal{C}(c, F-))$: cones from $c$ to $F$.

In enriched category theory, the relevant functor category is a
$\mathcal{V}$-functor category, and the ``constant diagram'' functor $\Delta c$
should be a $\mathcal{V}$-functor. But more general shapes arise naturally: we
may want to weight the limit by a functor $W : \mathcal{J}^{\mathrm{op}} \to
\mathcal{V}$ that assigns different ``multiplicities'' or ``shapes'' to each
object of $\mathcal{J}$.

A \term{$W$-weighted limit} of $F : \mathcal{J} \to \mathcal{C}$ is an object
$\{W, F\} \in \mathcal{C}$ characterized by the enriched universal property:
\begin{equation*}
  \mathcal{C}(c, \{W, F\}) \;\cong\; [\mathcal{J}^{\mathrm{op}}, \mathcal{V}]_\mathcal{V}(W, \mathcal{C}(c, F-))
\end{equation*}
as an isomorphism in $\mathcal{V}$, $\mathcal{V}$-natural in $c$.

\subsection{Expanding: Ordinary Limits as Special Cases}

Setting $\mathcal{V} = \mathbf{Set}$ and $W = \Delta 1$ (the constant weight
sending every object to the one-element set): the weighted limit $\{W, F\}$
recovers the ordinary limit $\lim F$, since $[\mathcal{J}, \mathbf{Set}](\Delta 1,
\mathcal{C}(c, F-))$ is the set of natural transformations from the constant
functor to $\mathcal{C}(c, F-)$, i.e., the set of cones.

More generally, taking $W = \mathcal{J}(j, -)$ (the representable weight at $j$)
and applying the Yoneda lemma gives $\{W, F\} \cong F(j)$: the limit weighted
by a representable recovers evaluation at that object.

\begin{example}{Cotensor as Weighted Limit}
For $\mathcal{V} = \mathbf{Ab}$ and $\mathcal{J} = \{*\}$ (one object), a
weight $W : \{*\}^{\mathrm{op}} \to \mathbf{Ab}$ is just an abelian group $A$,
and a functor $F : \{*\} \to \mathcal{C}$ is just an object $M \in \mathcal{C}$.
The weighted limit $\{A, M\}$ is the \emph{cotensor} of $M$ by $A$: the object
representing the functor $\mathcal{C}(-, M)^A$ in an appropriate sense. For
$\mathcal{C} = \mathbf{Ab}$, this is $\mathrm{Hom}(A, M)$.
\end{example}

\begin{example}{The Homotopy Limit}
In the $(\infty,1)$-categorical setting, homotopy limits are weighted limits
where the weight encodes simplicial resolution data. The classical Bousfield--Kan
formula for homotopy limits is a weighted colimit formula with the weight
coming from the nerve of the indexing category.
\end{example}

\subsection{Formalizing: The Definition and its Dual}

\begin{definition}
Let $\mathcal{C}$ be a $\mathcal{V}$-category, $\mathcal{J}$ a small
$\mathcal{V}$-category, $W : \mathcal{J}^{\mathrm{op}} \to \mathcal{V}$ a
$\mathcal{V}$-functor (the \term{weight}), and $F : \mathcal{J} \to \mathcal{C}$
a $\mathcal{V}$-functor. The \term{$W$-weighted limit} $\{W, F\}$, when it
exists, is an object of $\mathcal{C}$ with a $\mathcal{V}$-natural isomorphism
\begin{equation*}
  \mathcal{C}(c, \{W, F\})
  \;\cong\;
  [\mathcal{J}^{\mathrm{op}}, \mathcal{V}]_\mathcal{V}(W,\; \mathcal{C}(c, F(-))).
\end{equation*}

Dually, the \term{$W$-weighted colimit} $W \star F$, when it exists, is an
object with
\begin{equation*}
  \mathcal{C}(W \star F,\; c)
  \;\cong\;
  [\mathcal{J}, \mathcal{V}]_\mathcal{V}(W,\; \mathcal{C}(F(-), c)).
\end{equation*}
\end{definition}

\subsection{Reorganizing: The Yoneda Reduction}

The reason ordinary limits are a special case is the Yoneda lemma: for $W =
\Delta I$ (constant weight at the unit object $I \in \mathcal{V}$),
\begin{equation*}
  [\mathcal{J}^{\mathrm{op}}, \mathcal{V}]_\mathcal{V}(\Delta I,\; \mathcal{C}(c, F(-)))
  \;\cong\;
  \int_j \mathcal{V}(I, \mathcal{C}(c, F(j)))
  \;\cong\;
  \int_j \mathcal{C}(c, F(j)),
\end{equation*}
the last step using $\mathcal{V}(I, -) \cong (-)$ (since $I$ is the unit).
For $\mathcal{V} = \mathbf{Set}$ this is the set of natural cones, i.e., ordinary
limit cones.

\begin{howtosay}
\begin{itemize}
  \item $\{W, F\}$ is pronounced ``$W$ hom $F$'' or ``the $W$-weighted limit
        of $F$.'' Some authors write $\mathrm{lim}^W F$ or $[W, F]$.
  \item $W \star F$ is the \emph{coend} or weighted colimit; also written
        $\mathrm{colim}^W F$ or $W \otimes F$ in older texts.
  \item The $\Delta I$-weighted limit is the ordinary conical limit; the
        $\Delta I$-weighted colimit is the ordinary conical colimit.
  \item Weighted limits subsume: ordinary limits, ends, cotensors, homotopy
        limits, and mapping spaces.
\end{itemize}
\end{howtosay}


% ═══════════════════════════════════════════════════════════════════════════
\section{The Change of Base}
\label{sec:change-of-base}
% ═══════════════════════════════════════════════════════════════════════════

\subsection{Informally: Moving Between Enrichments}

If $\Phi : \mathcal{V} \to \mathcal{W}$ is a lax monoidal functor (Chapter~19),
we can use $\Phi$ to turn any $\mathcal{V}$-category into a $\mathcal{W}$-category
by applying $\Phi$ to all hom-objects. This is the \term{change of base}.

\subsection{Expanding: The Construction}

Let $\mathcal{C}$ be a $\mathcal{V}$-category. The \term{base change} along
$\Phi : \mathcal{V} \to \mathcal{W}$ is the $\mathcal{W}$-category
$\mathcal{C}_\Phi$ with:
\begin{itemize}
  \item The same objects as $\mathcal{C}$.
  \item Hom-objects $\mathcal{C}_\Phi(a, b) = \Phi(\mathcal{C}(a, b))$ in $\mathcal{W}$.
  \item Composition $\mathcal{C}_\Phi(b,c) \otimes_\mathcal{W} \mathcal{C}_\Phi(a,b)
        \to \mathcal{C}_\Phi(a,c)$ defined as
        \begin{equation*}
          \Phi(\mathcal{C}(b,c)) \otimes_\mathcal{W} \Phi(\mathcal{C}(a,b))
          \xrightarrow{\phi_{-,-}} \Phi(\mathcal{C}(b,c) \otimes_\mathcal{V}
          \mathcal{C}(a,b)) \xrightarrow{\Phi(\circ_\mathcal{C})} \Phi(\mathcal{C}(a,c)),
        \end{equation*}
        using the lax monoidal comparison $\phi$ of $\Phi$.
  \item Identities $j_a^\mathcal{W} : I_\mathcal{W} \xrightarrow{\phi_0}
        \Phi(I_\mathcal{V}) \xrightarrow{\Phi(j_a^\mathcal{C})} \Phi(\mathcal{C}(a,a))$.
\end{itemize}

The coherence conditions for $\mathcal{V}$-categories (associativity and unit
laws for $\circ_\mathcal{C}$) together with the coherence of the lax monoidal
functor $\Phi$ guarantee that $\mathcal{C}_\Phi$ is a valid $\mathcal{W}$-category.

\subsection{Formalizing: Functoriality of Base Change}

Base change is functorial: if $F : \mathcal{C} \to \mathcal{D}$ is a
$\mathcal{V}$-functor, then applying $\Phi$ to the hom-object maps $F_{a,b} :
\mathcal{C}(a,b) \to \mathcal{D}(Fa, Fb)$ gives $\mathcal{W}$-functor maps
$\Phi(F_{a,b}) : \mathcal{C}_\Phi(a,b) \to \mathcal{D}_\Phi(Fa,Fb)$, yielding
a $\mathcal{W}$-functor $F_\Phi : \mathcal{C}_\Phi \to \mathcal{D}_\Phi$.

\begin{proposition}
The change-of-base construction defines a 2-functor
$(-)_\Phi : \mathcal{V}\text{-}\mathbf{Cat} \to \mathcal{W}\text{-}\mathbf{Cat}$.
\end{proposition}

\subsection{Reorganizing: Examples of Base Change}

\begin{example}{Forgetting the Enrichment}
Let $\Phi = \mathcal{V}(I,-) : \mathcal{V} \to \mathbf{Set}$. This is lax
monoidal (it sends $\otimes$ to $\times$ via the composition map). For any
$\mathcal{V}$-category $\mathcal{C}$, the base change $\mathcal{C}_\Phi$ is the
\term{underlying ordinary category} of $\mathcal{C}$: its hom-sets are
$\mathcal{V}(I, \mathcal{C}(a,b))$.
\end{example}

\begin{example}{Complexes to Graded}
The cohomology functor $H : \mathbf{Ch}(\mathbf{Ab}) \to \mathbf{GrAb}$ (from
chain complexes to graded abelian groups) is lax monoidal with respect to
tensor products, and induces a base change from $\mathbf{Ch}(\mathbf{Ab})$-enriched
to $\mathbf{GrAb}$-enriched categories.
\end{example}


% ═══════════════════════════════════════════════════════════════════════════
\section{Self-Enrichment of Closed Symmetric Monoidal Categories}
\label{sec:self-enrichment}
% ═══════════════════════════════════════════════════════════════════════════

\subsection{Informally: A Category Enriching Over Itself}

The most important source of enriched categories: a \term{closed symmetric
monoidal category} $(\mathcal{V}, \otimes, I, [-,-])$ enriches over itself.
Here $[A, B]$ is the \emph{internal hom} object, characterized by:
\begin{equation*}
  \mathcal{V}(C \otimes A, B) \;\cong\; \mathcal{V}(C, [A,B]).
\end{equation*}
This is the \term{closed structure}: $(- \otimes A)$ has a right adjoint $[A,-]$.

\begin{dialog}
Why does closedness give enrichment?

To enrich $\mathcal{V}$ over itself, we need hom-objects in $\mathcal{V}$.
The internal hom $[A, B]$ is exactly the right object to use. The composition
law $[B,C] \otimes [A,B] \to [A,C]$ (``evaluation followed by application'')
is the standard map in a closed category, and the identity $I \to [A,A]$
comes from the unit isomorphism.
\end{dialog}

\subsection{Expanding: The Self-Enrichment}

Let $(\mathcal{V}, \otimes, I)$ be a closed symmetric monoidal category with
internal hom $[-, -]$. Define a $\mathcal{V}$-category (also called $\mathcal{V}$):
\begin{itemize}
  \item Objects: the objects of $\mathcal{V}$.
  \item $\mathcal{V}$-hom-objects: $\mathcal{V}(A, B) = [A, B] \in \mathcal{V}$.
  \item Composition: the standard morphism $[B,C] \otimes [A,B] \to [A,C]$
        obtained by currying the evaluation composite
        $[B,C] \otimes [A,B] \otimes A \to [B,C] \otimes B \to C$.
  \item Identity: the morphism $I \to [A,A]$ corresponding (under the adjunction)
        to the identity $A \otimes I \cong A$.
\end{itemize}
These satisfy the $\mathcal{V}$-category axioms by the coherence of the closed
structure.

\subsection{Formalizing: Examples}

\begin{example}{$\mathbf{Ab}$-enrichment}
$(\mathbf{Ab}, \otimes_\mathbb{Z}, \mathbb{Z})$ is closed symmetric monoidal
with internal hom $[A, B] = \mathrm{Hom}(A, B)$ (the group of homomorphisms).
The self-enrichment says: $\mathbf{Ab}$ is enriched over itself, with hom-objects
$\mathrm{Hom}(A, B)$ being abelian groups. This is the standard
$\mathbf{Ab}$-enrichment of the category of abelian groups.
\end{example}

\begin{example}{$\mathbf{Vect}_k$-enrichment}
$(\mathbf{Vect}_k, \otimes_k, k)$ is closed with $[V, W] = \mathrm{Hom}_k(V, W)$.
The self-enrichment gives $\mathbf{Vect}_k$ as a $\mathbf{Vect}_k$-enriched
category, with hom-objects being vector spaces of linear maps.
\end{example}

\begin{example}{$\mathbf{sSet}$-enrichment}
The category of simplicial sets $\mathbf{sSet} = [\Delta^{\mathrm{op}}, \mathbf{Set}]$
is closed monoidal with the cartesian product. The internal hom is the
simplicial mapping space $\mathbf{Map}(X, Y)_n = \mathbf{sSet}(X \times \Delta^n, Y)$.
Self-enrichment gives $\mathbf{sSet}$ as an $\mathbf{sSet}$-enriched category,
which is the starting point for simplicial homotopy theory.
\end{example}

\subsection{Reorganizing: Tensors, Cotensors, and Completeness}

A $\mathcal{V}$-category $\mathcal{C}$ is \term{$\mathcal{V}$-complete} if all
weighted limits $\{W, F\}$ exist. It is \term{$\mathcal{V}$-tensored} if for
each $V \in \mathcal{V}$ and $c \in \mathcal{C}$, the tensor $V \otimes c$
exists, characterized by $\mathcal{C}(V \otimes c, d) \cong \mathcal{V}(V,
\mathcal{C}(c, d))$. Dually for cotensors $[V, c]$.

When $\mathcal{C}$ is both tensored and cotensored over $\mathcal{V}$, all
weighted limits and colimits reduce to ordinary conical limits plus the
(co)tensor action. This is a powerful computational tool.

\begin{theorem}
A $\mathcal{V}$-category $\mathcal{C}$ that is $\mathcal{V}$-tensored and
$\mathcal{V}$-cotensored has all weighted limits iff it has all conical limits:
\begin{equation*}
  \{W, F\} \;\cong\; \lim_{j \in \mathcal{J}} [W(j), F(j)],
\end{equation*}
where $[W(j), F(j)]$ denotes the cotensor of $F(j)$ by $W(j)$.
\end{theorem}


% ═══════════════════════════════════════════════════════════════════════════
\section*{Exercises}
\addcontentsline{toc}{section}{Exercises}
% ═══════════════════════════════════════════════════════════════════════════

\begin{exercisesbox}
\begin{qa}
\qaitem{State the enriched Yoneda lemma.}{$[\mathcal{C},\mathcal{V}]_\mathcal{V}(\mathcal{C}(a,-), F) \cong F(a)$ in $\mathcal{V}$, $\mathcal{V}$-naturally in $a$.}
\qaitem{What is the $\mathcal{V}$-object of $\mathcal{V}$-natural transformations?}{$\int_c \mathcal{V}(G(c), H(c))$, an end in $\mathcal{V}$.}
\qaitem{In the enriched Yoneda proof, what is $\phi(\theta)$?}{$\phi(\theta) = \theta_a \circ j_a : I \to F(a)$, where $j_a : I \to \mathcal{C}(a,a)$ is the enriched identity.}
\qaitem{What is a $\mathcal{V}$-adjunction?}{A $\mathcal{V}$-natural isomorphism $\mathcal{D}(FA,B) \cong \mathcal{C}(A,GB)$ in $\mathcal{V}$.}
\qaitem{Does a $\mathcal{V}$-adjunction imply an ordinary adjunction?}{Yes: apply $\mathcal{V}(I,-)$ to get a bijection of underlying hom-sets.}
\qaitem{Can an ordinary adjunction between $\mathcal{V}$-functors fail to be a $\mathcal{V}$-adjunction?}{Yes, in general. But in tensored/cotensored $\mathcal{V}$-categories, ordinary adjoints between $\mathcal{V}$-functors lift to $\mathcal{V}$-adjunctions.}
\qaitem{What is a $W$-weighted limit?}{An object $\{W,F\} \in \mathcal{C}$ with $\mathcal{C}(c,\{W,F\}) \cong [\mathcal{J}^{\mathrm{op}},\mathcal{V}]_\mathcal{V}(W, \mathcal{C}(c,F-))$ $\mathcal{V}$-naturally in $c$.}
\qaitem{What ordinary limit corresponds to $W = \Delta I$?}{The ordinary conical limit $\lim F$.}
\qaitem{What weighted limit corresponds to $W = \mathcal{J}(j,-)$?}{Evaluation $F(j)$, by the enriched Yoneda lemma.}
\qaitem{What is the cotensor $[A, M]$ as a weighted limit?}{The $A$-weighted limit of the single-object diagram at $M$, for $A \in \mathcal{V}$ and $M \in \mathcal{C}$.}
\qaitem{What is a closed symmetric monoidal category?}{One where $(-\otimes A)$ has a right adjoint $[A,-]$ (the internal hom), and the tensor is symmetric.}
\qaitem{What is self-enrichment?}{Using $[A,B]$ as the hom-object of $A,B$ in the $\mathcal{V}$-category structure on $\mathcal{V}$ itself.}
\qaitem{What is the internal hom in $\mathbf{Ab}$?}{$[A,B] = \mathrm{Hom}(A,B)$, the abelian group of homomorphisms.}
\qaitem{What is the change of base for enriched categories?}{Given a lax monoidal $\Phi : \mathcal{V} \to \mathcal{W}$, a $\mathcal{V}$-category $\mathcal{C}$ becomes a $\mathcal{W}$-category $\mathcal{C}_\Phi$ with hom-objects $\Phi(\mathcal{C}(a,b))$.}
\qaitem{What lax monoidal functor gives the underlying ordinary category?}{$\mathcal{V}(I,-) : \mathcal{V} \to \mathbf{Set}$, sending each hom-object to its underlying set of ``generalized elements.'' }
\qaitem{What is a $\mathcal{V}$-tensored category?}{One where $V \otimes c$ exists with $\mathcal{C}(V\otimes c, d) \cong \mathcal{V}(V, \mathcal{C}(c,d))$.}
\qaitem{In a tensored and cotensored $\mathcal{V}$-category, how do weighted limits reduce?}{$\{W,F\} \cong \lim_j [W(j), F(j)]$, a conical limit of cotensors.}
\qaitem{What is the enriched Yoneda embedding?}{$Y : \mathcal{C} \to [\mathcal{C}^{\mathrm{op}},\mathcal{V}]_\mathcal{V}$, $Y(a) = \mathcal{C}(-,a)$; it is fully faithful as a $\mathcal{V}$-functor.}
\qaitem{What does $\mathcal{V}$-naturality require beyond ordinary naturality?}{The naturality squares must commute in $\mathcal{V}$ (as morphisms in $\mathcal{V}$), not just in $\mathbf{Set}$.}
\qaitem{What is the simplicial mapping space in $\mathbf{sSet}$?}{$\mathbf{Map}(X,Y)_n = \mathbf{sSet}(X \times \Delta^n, Y)$, the internal hom in the closed monoidal category $\mathbf{sSet}$.}
\end{qa}
\end{exercisesbox}


% ═══════════════════════════════════════════════════════════════════════════
\section*{Chapter Summary}
\addcontentsline{toc}{section}{Chapter Summary}
% ═══════════════════════════════════════════════════════════════════════════

\begin{summarybox}
\textbf{Enriched Yoneda.} For a $\mathcal{V}$-functor $F : \mathcal{C} \to
\mathcal{V}$, $[\mathcal{C},\mathcal{V}]_\mathcal{V}(\mathcal{C}(a,-), F)
\cong F(a)$ in $\mathcal{V}$. The proof uses the enriched identity morphism
and the $\mathcal{V}$-functor action, with mutual inverse $\psi(x)_c = F_{a,c}
\circ (x \otimes \mathrm{id})$.

\textbf{$\mathcal{V}$-adjunctions.} The ordinary hom-set bijection is replaced
by a $\mathcal{V}$-natural isomorphism $\mathcal{D}(FA,B) \cong \mathcal{C}(A,GB)$
in $\mathcal{V}$. Every $\mathcal{V}$-adjunction has an underlying ordinary
adjunction.

\textbf{Weighted limits.} The $W$-weighted limit $\{W,F\}$ satisfies
$\mathcal{C}(c, \{W,F\}) \cong [\mathcal{J}^{\mathrm{op}},\mathcal{V}]_\mathcal{V}(W,
\mathcal{C}(c, F-))$. Ordinary limits are $\Delta I$-weighted. In tensored and
cotensored categories, $\{W,F\} \cong \lim_j [W(j), F(j)]$.

\textbf{Change of base.} A lax monoidal functor $\Phi : \mathcal{V} \to \mathcal{W}$
induces $(-)_\Phi : \mathcal{V}\text{-}\mathbf{Cat} \to \mathcal{W}\text{-}\mathbf{Cat}$
by applying $\Phi$ to all hom-objects. The special case $\mathcal{V}(I,-)$ gives
the underlying ordinary category.

\textbf{Self-enrichment.} A closed symmetric monoidal category $(\mathcal{V},
\otimes, I, [-,-])$ enriches over itself using $[A,B]$ as the hom-object.
Examples: $\mathbf{Ab}$, $\mathbf{Vect}_k$, $\mathbf{sSet}$.
\end{summarybox}


% ═══════════════════════════════════════════════════════════════════════════
\section*{Formal Restatement}
\addcontentsline{toc}{section}{Formal Restatement}
% ═══════════════════════════════════════════════════════════════════════════

\begin{formalrestatement}
\textbf{Enriched Yoneda.} For $\mathcal{C}$ a $\mathcal{V}$-category and
$F : \mathcal{C} \to \mathcal{V}$ a $\mathcal{V}$-functor:
\begin{equation*}
  [\mathcal{C},\mathcal{V}]_\mathcal{V}(\mathcal{C}(a,-), F) \;=\;
  \int_c \mathcal{V}(\mathcal{C}(a,c), F(c)) \;\cong\; F(a)
\end{equation*}
$\mathcal{V}$-naturally in $a$.

\medskip
\textbf{$\mathcal{V}$-adjunction.} $F \dashv_\mathcal{V} G$ iff there is a
$\mathcal{V}$-natural isomorphism $\mathcal{D}(FA,B) \cong \mathcal{C}(A,GB)$
in $\mathcal{V}$. Underlying ordinary adjunction obtained by applying
$\mathcal{V}(I,-)$.

\medskip
\textbf{Weighted limit.}
\begin{equation*}
  \mathcal{C}(c, \{W,F\}) \;\cong\;
  [\mathcal{J}^{\mathrm{op}},\mathcal{V}]_\mathcal{V}(W, \mathcal{C}(c,F-)).
\end{equation*}
Ordinary limit: $W = \Delta I$. In tensored/cotensored $\mathcal{C}$: $\{W,F\}
\cong \lim_j [W(j), F(j)]$.

\medskip
\textbf{Change of base.} Lax monoidal $\Phi : \mathcal{V} \to \mathcal{W}$
induces $\mathcal{C}_\Phi$ with $\mathcal{C}_\Phi(a,b) = \Phi(\mathcal{C}(a,b))$,
composites via $\phi$-comparisons.

\medskip
\textbf{Self-enrichment.} Closed symmetric monoidal $(\mathcal{V},\otimes,I,[-,-])$
is a $\mathcal{V}$-category with $\mathcal{V}(A,B) = [A,B]$, composition via
the evaluation map $[B,C]\otimes[A,B]\otimes A \to C$.

\medskip
\noindent\textit{Chapter~21 returns to monads via Beck's monadicity theorem,
which characterizes when a functor is equivalent to a category of algebras
for a monad.}
\end{formalrestatement}
